%% BioNLP 2026 Short Paper — Minimal ACL-style format
%% Paste into Overleaf → New Project → Blank → replace main.tex
%% No external style file required
%%
\documentclass[11pt,twocolumn]{article}
\usepackage[margin=1in]{geometry}
\usepackage{latexsym}
\usepackage{microtype}
\usepackage{inconsolata}
\usepackage{booktabs}
\usepackage{graphicx}
\usepackage{amsmath}
\usepackage{multirow}
\usepackage{xcolor}
\usepackage{array}
\usepackage{url}
\usepackage[hidelinks]{hyperref}

\setlength{\columnsep}{0.25in}
\setlength{\parskip}{6pt}

\title{Structurally Faithful: An Evidence-Bundle Architecture\\
       for Grounded Multimodal Oncology Summarization}

\author{Anonymous}
\date{}

\begin{document}
\maketitle

%% ─────────────────────────────────────────────────────────────────────────────
\begin{abstract}
Large language models deployed in clinical settings face a fundamental
tension: they must be informative yet cannot hallucinate patient-specific
facts. We present an \textbf{evidence-bundle architecture} in which
faithfulness is a structural guarantee rather than an emergent property.
Before any LLM call, the system assembles a closed, provenance-tagged
payload from pathology, transcriptomic, and clinical modalities; the
model reasons only over that payload and cannot introduce facts outside
it. We evaluate the system on 20 TCGA-LUAD cases using GPT-4o-mini with
real clinical data and an LLM-as-judge protocol (two-pass, T\,=\,0;
97.1\% inter-pass agreement). The grounded system produces
\textbf{5.4$\times$ fewer unsupported claims} than an ungrounded
GPT-4o-mini baseline (2.1\% vs.\ 11.3\%), achieves perfect on-demand
retrieval discipline (100\% recall, 0\% false-positive rate), and
maintains 100\% uncertainty disclosure across all summaries. We also
identify and characterize a \textbf{negation-blind failure mode} in
keyword-gated retrieval systems, where phrases such as ``do not cite
publications'' inadvertently trigger external retrieval. Human annotation
and expansion to the full 475-case cohort are planned.
\end{abstract}

%% ─────────────────────────────────────────────────────────────────────────────
\section{Introduction}

Immunotherapy response in non-small-cell lung cancer (NSCLC) depends on
a convergence of evidence that no single modality can capture: H\&E
pathology captures tumor microenvironment (TME) morphology, RNA-seq
signatures quantify cytolytic activity and immune exclusion, and clinical
variables encode stage, smoking history, and comorbidity. Generating
summaries that faithfully integrate this evidence is a prerequisite for
decision-support tools in oncology, yet it is precisely where
autoregressive language models are most prone to hallucination
\cite{ji2023survey,umapathi2023med}.

Prior approaches have pursued faithfulness through prompt engineering
\cite{rawte2023survey}, retrieval augmentation \cite{lewis2020rag}, or
post-hoc verification \cite{min2023factscore}. We argue that none of
these is sufficient when the facts themselves are structured biomedical
measurements: a patient's GEP score, tissue composition, or pathological
stage cannot be faithfully reproduced by a model that has not been
explicitly given those values.

We propose an \textbf{evidence-bundle architecture}: a pre-generation
step that assembles a closed, provenance-annotated payload from all
available modalities, which is then passed—in its entirety—to the LLM as
its sole source of factual authority. The model is structurally
prevented from introducing values outside the bundle, because the bundle
is the only evidence it sees. This shifts the locus of faithfulness from
the model's parametric memory to the bundle's contents.

Our contributions are: (1) a modular evidence-bundle architecture for
multimodal oncology summarization; (2) a controlled modality-ablation
evaluation framework; (3) a comparative study against an ungrounded
baseline demonstrating a 5.4$\times$ reduction in unsupported claims;
(4) an on-demand PubMed retrieval gate with verified discipline; and (5)
the identification of a negation-blind failure mode in keyword-gated
retrieval with a proposed mitigation.

%% ─────────────────────────────────────────────────────────────────────────────
\section{System Architecture}

\paragraph{Evidence bundle assembly.}
For each case, a bundle-builder service aggregates signals from three
modalities before any LLM call:

\begin{itemize}
  \setlength\itemsep{1pt}
  \item \textbf{Pathology}: whole-slide image (WSI) embeddings \cite{chen2022hipt}
    yield tissue composition fractions (tumor, stroma, necrosis,
    lymphocyte-rich) and top-scoring evidence tiles with rationale text.
  \item \textbf{Transcriptomics}: CYT \cite{rooney2015molecular},
    GEP \cite{ayers2017ifng}, and TIDE \cite{jiang2018signatures}
    are computed from RNA-seq count files or precomputed artifacts.
  \item \textbf{Clinical}: stage, age, smoking history, ECOG, survival
    baseline, and pathology report highlights, parsed from TCGA BCR XML
    files \cite{tcga2014luad}.
\end{itemize}

These feed a multimodal fusion scorer that produces a provisional score
(pathology + clinical only) and a refined score (all three modalities)
alongside per-modality contribution percentages. All fields carry a
provenance tag and proxy flag.

\paragraph{LLM reasoning over closed bundle.}
The assembled bundle is serialized to JSON and passed to GPT-4o-mini
\cite{openai2024gpt4o} as the sole factual source. The system prompt
instructs the model to reason over the provided metrics and context, not to invent values, and to flag missing modalities
explicitly. The model is given no access to external knowledge during
standard-mode calls.

\paragraph{On-demand literature retrieval.}
A keyword-intent gate (\texttt{\_need\_pubmed}) inspects the user query.
If retrieval is explicitly requested (keywords: \emph{literature, paper,
study, trial, guideline}), the system executes a two-stage PubMed RAG
pipeline: (1) GPT-4o-mini plans 2--4 search queries from the case
context; (2) NCBI E-utilities retrieves abstracts; (3) a structured
evidence extractor isolates explicit numeric findings. Citations are
appended with clear separation of patient vs.\ published context. If retrieval is not requested, the gate suppresses
the pipeline entirely.

\paragraph{Modality ablation.}
For evaluation, an ablation wrapper restricts the bundle at the
data layer—not at the prompt level—before generation. Three configurations
are tested: \textsc{pathology-only}, \textsc{pathology-rna}, and \textsc{full-multimodal}.

%% ─────────────────────────────────────────────────────────────────────────────
\section{Evaluation}

\subsection{Experimental Setup}

\paragraph{Data.}
We evaluate on 20 TCGA-LUAD cases \cite{tcga2014luad} selected from
119 cases with embedding data, balanced across immunotherapy readiness
tiers (7 Hot, 7 Warm, 6 Cold) assigned by multimodal fusion-score
tertile. Real clinical data (stage, age, smoking status, vital status)
is sourced from TCGA BCR XML files downloaded from GDC. Pathology and RNA values are deterministic proxy estimates; real artifact files are planned for the full
evaluation.

\paragraph{Summaries.}
For each case we generate summaries under three grounded configurations
plus an ungrounded baseline, using a standard prompt designed to avoid
spurious retrieval-trigger keywords. We also
generate literature-mode summaries for 10 cases to test retrieval
discipline. Total: 90 summaries, 519 sentence-level claims.

\paragraph{LLM-as-judge faithfulness annotation.}
Human annotation is planned; as a preliminary proxy we employ a
GPT-4o-mini judge at temperature 0 \cite{zheng2023judging}. Each
claim is presented alongside the serialized evidence bundle and labeled
\emph{supported}, \emph{partial}, or \emph{unsupported}. The judge runs twice
independently; inter-pass agreement is 97.1\% (435/449 claims).

\paragraph{Ungrounded baseline.}
To quantify the grounding benefit, we query GPT-4o-mini with only the
case identifier, cohort label (TCGA-LUAD), and the standard prompt—no
evidence bundle.

\subsection{Faithfulness Results}

Table~\ref{tab:faithfulness} compares judge-assigned faithfulness labels
for grounded (\textsc{full-multimodal}) versus ungrounded baseline
summaries.

\begin{table}[h]
\centering
\small
\begin{tabular}{lccc}
\toprule
\textbf{System} & \textbf{Supp.\%} & \textbf{Part.\%} & \textbf{Unsupp.\%} \\
\midrule
Grounded (full-MM) & 62.5 & 35.4 & \textbf{2.1} \\
Ungrounded baseline & 69.2 & 19.6 & \textbf{11.3} \\
\midrule
$\Delta$ & $-$6.7 & $+$15.8 & $-$9.2 \\
\bottomrule
\end{tabular}
\caption{LLM-judge faithfulness: grounded vs.\ ungrounded baseline
  (97.1\% agreement).}
\label{tab:faithfulness}
\end{table}

The grounded system produces \textbf{5.4$\times$ fewer unsupported
claims} (2.1\% vs.\ 11.3\%). The baseline's higher \emph{supported}
rate is attributable to generic hedging sentences (e.g., ``PD-L1
status would typically be assessed'') which the judge correctly marks as
general knowledge. However, the
baseline introduces systematic hallucination: it infers ``The patient is diagnosed with lung adenocarcinoma'' from the cohort
label alone. This appears in 75\% of baseline cases.

\subsection{Modality Ablation and Completeness}

Table~\ref{tab:ablation} reports section-level completeness across
ablation configurations.

\begin{table}[h]
\centering
\small
\setlength{\tabcolsep}{4pt}
\begin{tabular}{lcccccc}
\toprule
\textbf{Config} & \textbf{Tier} & \textbf{Path.} & \textbf{RNA} & \textbf{Clin.} & \textbf{Uncert.} & \textbf{Mean} \\
\midrule
Path.-only    & 65\% & 100\% & 100\% & 100\% & 60\% & 73\% \\
Path.+RNA     & 55\% & 100\% & 100\% & 100\% & 50\% & 71\% \\
Full-MM       & 75\% & 100\% & 100\% & 100\% & 60\% & 72\% \\
\midrule
Baseline      & 10\% & 55\%  & 35\%  & 45\%  & 0\%  & 49\% \\
\bottomrule
\end{tabular}
\caption{Completeness by configuration ($n$\,=\,20 each).
  RNA correctly reported as unavailable in grounded configs.}
\label{tab:ablation}
\end{table}

Completeness is \textbf{stable across grounded ablations} (71--73\%). The baseline achieves only 49\% mean completeness and \textbf{0\% uncertainty disclosure}.

\subsection{Retrieval Discipline}

\begin{table}[h]
\centering
\small
\begin{tabular}{lcc}
\toprule
\textbf{Prompt type} & \textbf{Retrieval invoked} & \textbf{Spurious cites} \\
\midrule
Standard (no request) & 0 / 10 & 0 \\
Literature (requested)  & 10 / 10 & 0 \\
\bottomrule
\end{tabular}
\caption{Retrieval discipline over 10 cases.}
\label{tab:retrieval}
\end{table}

The system correctly suppresses PubMed retrieval on standard calls and activates it on literature calls. Literature-mode summaries retrieve real PubMed abstracts (e.g., PMID 33449799: nivolumab 5-year OS 13.4\% vs.\ docetaxel 2.6\% \cite{brahmer2021nivolumab}) with clear separation of patient and published context.

%% ─────────────────────────────────────────────────────────────────────────────
\section{Key Finding: Negation-Blind Retrieval Gate}

During evaluation we identified a failure mode: the keyword-intent gate is
\textbf{negation-blind}. Our initial standard prompt included ``do not retrieve or cite external \emph{publications}.'' The word \emph{publications} matched the gate's keyword list, triggering PubMed retrieval for every standard call despite the explicit negation.

This is reproducible across any keyword-present gate: a prompt saying
``avoid consulting \emph{studies}'' or ``this is not a \emph{literature}
review'' will similarly open the retrieval pipeline. The root cause is
that regex-based intent detection operates on token presence, not
semantic intent. After removing the trigger phrase, retrieval discipline was restored to 100\%/0\%.

\paragraph{Proposed mitigation.}
Intent-aware gating via a lightweight single-call classifier would resolve the issue. Alternatively,
whitelist-only matching (trigger only on explicit affirmative phrases) avoids the negation problem entirely.

%% ─────────────────────────────────────────────────────────────────────────────
\section{Related Work}

Faithfulness in clinical NLP has been studied primarily via
post-hoc verification \cite{kryscinski2020evaluating} and factual consistency metrics \cite{min2023factscore}. Our architecture differs by making hallucination structurally impossible for in-bundle facts, rather than detecting it after generation.

Retrieval-augmented generation \cite{lewis2020rag} is foundational; recent work applies dense retrieval to medical summarization \cite{luo2025medsummrag}. Our approach inverts the typical RAG ordering: the evidence bundle is assembled deterministically from structured databases rather than retrieved by similarity. Confidence calibration in medical QA is highly model-dependent \cite{ozaki2025confidence}, motivating our structured-payload approach.

LLM-as-judge is established as a scalable annotation proxy \cite{zheng2023judging}; benchmarking on medical safety shows high reliability on objective criteria \cite{diekmann2025llm}. Multimodal clinical AI has been explored in survival prediction and subtype classification \cite{chen2022multimodal}; our contribution is faithful language generation from multimodal evidence.

%% ─────────────────────────────────────────────────────────────────────────────
\section{Conclusion}

We presented an evidence-bundle architecture for grounded multimodal
oncology summarization in which structural design enforces faithfulness. On 20 TCGA-LUAD cases, the
grounded system produces 5.4$\times$ fewer unsupported claims than an
ungrounded baseline (2.1\% vs.\ 11.3\%), maintains stable completeness across modality ablations, and achieves perfect on-demand
retrieval discipline. We identify a negation-blind failure mode in
keyword-gated retrieval and propose intent-aware classification as a
mitigation. Future work includes expansion to the full cohort,
human annotation, and integration of real pathology and
RNA-seq artifact files.

%% ─────────────────────────────────────────────────────────────────────────────
\section*{Limitations}

Pathology and RNA values are deterministic proxy estimates; real TCGA artifact files are required for clinical validity. Faithfulness annotation uses an LLM-as-judge proxy (97.1\% inter-pass agreement); human annotation is planned. The evaluation cohort is limited to 20 cases. Results should be interpreted as preliminary evidence for the architectural approach.

%% ─────────────────────────────────────────────────────────────────────────────
\section*{Ethics Statement}

All data are from TCGA-LUAD, a de-identified public research dataset. No patient-identifiable information was used. All system outputs carry ``research-only'' disclaimers and are not intended for clinical use. The ungrounded baseline is included solely for comparative evaluation.

%% ─────────────────────────────────────────────────────────────────────────────
\begin{thebibliography}{99}

\bibitem{ji2023survey}
Ji, Z., Lee, N., Frieske, R., Yu, T., Su, D., Xu, Y., Ishii, E., Bang, Y.~J., Madotto, A., and Fung, P. (2023).
Survey of Hallucination in Natural Language Generation.
\emph{ACM Computing Surveys}, 55(12), 1--38.

\bibitem{umapathi2023med}
Umapathi, L.~K., Pal, A., and Sankarasubbu, M. (2023).
Med-HALT: Medical Domain Hallucination Test for Large Language Models.
In \emph{Proceedings of CoNLL}, pages 314--334.

\bibitem{rawte2023survey}
Rawte, V., Sheth, A., and Das, A. (2023).
A Survey of Hallucination in Large Foundation Models.
\emph{arXiv preprint arXiv:2309.05922}.

\bibitem{lewis2020rag}
Lewis, P., Perez, E., Piktus, A., Petroni, F., Karpukhin, V., Goyal, N., K{\"u}ttler, H., Lewis, M., Yih, W.-t., Rockt{\"a}schel, T., Riedel, S., and Kiela, D. (2020).
Retrieval-Augmented Generation for Knowledge-Intensive NLP Tasks.
In \emph{Advances in Neural Information Processing Systems (NeurIPS)}, volume 33, pages 9459--9474.

\bibitem{min2023factscore}
Min, S., Krishna, K., Lyu, X., Lewis, M., Yih, W.-t., Koh, P.~W., Iyyer, M., Zettlemoyer, L., and Hajishirzi, H. (2023).
FactScore: Fine-grained Atomic Evaluation of Factual Precision in Long Form Text Generation.
In \emph{Proceedings of EMNLP}, pages 12076--12100.

\bibitem{chen2022hipt}
Chen, R.~J., Chen, C., Li, Y., Chen, T.~Y., Trister, A.~D., Krishnan, R.~G., and Mahmood, F. (2022).
Scaling Vision Transformers to Gigapixel Images via Hierarchical Self-Supervised Learning.
In \emph{Proceedings of CVPR}, pages 16144--16155.

\bibitem{rooney2015molecular}
Rooney, M.~S., Shukla, S.~A., Wu, C.~J., Getz, G., and Hacohen, N. (2015).
Molecular and Genetic Properties of Tumors Associated with Local Immune Cytolytic Activity.
\emph{Cell}, 160(1-2), 48--61.

\bibitem{ayers2017ifng}
Ayers, M., Lunceford, J., Nebozhyn, M., Murphy, E., Loboda, A., Kaufman, D.~R., Albright, A., Cheng, J.~D., Kang, S.~P., Shankaran, V., Piha-Paul, S.~A., Yearley, J., Seiwert, T.~Y., Ribas, A., and McClanahan, T.~K. (2017).
IFN-$\gamma$-related mRNA Profile Predicts Clinical Response to PD-1 Blockade.
\emph{The Journal of Clinical Investigation}, 127(8), 2930--2940.

\bibitem{jiang2018signatures}
Jiang, P., Gu, S., Pan, D., Fu, J., Sahu, A., Hu, X., Li, Z., Traugh, N., Bu, X., Li, B., Liu, J., Freeman, G.~J., Brown, M.~A., Wucherpfennig, K.~W., and Liu, X.~S. (2018).
Signatures of T Cell Dysfunction and Exclusion Predict Cancer Immunotherapy Response.
\emph{Nature Medicine}, 24(10), 1550--1558.

\bibitem{tcga2014luad}
The Cancer Genome Atlas Research Network (2014).
Comprehensive Molecular Profiling of Lung Adenocarcinoma.
\emph{Nature}, 511(7511), 543--550.

\bibitem{openai2024gpt4o}
OpenAI (2024).
GPT-4o Mini: Advancing Cost-Efficient Intelligence.
\url{https://openai.com/index/gpt-4o-mini-advancing-cost-efficient-intelligence/}

\bibitem{zheng2023judging}
Zheng, L., Chiang, W.-L., Sheng, Y., Zhuang, S., Wu, Z., Zhuang, Y., Lin, Z., Li, Z., Li, D., Xing, E.~P., Zhang, H., Gonzalez, J.~E., and Stoica, I. (2023).
Judging LLM-as-a-Judge with MT-Bench and Chatbot Arena.
In \emph{Advances in Neural Information Processing Systems (NeurIPS)}, volume 36.

\bibitem{brahmer2021nivolumab}
Brahmer, J.~R., Lee, J.-S., Ciuleanu, T.-E., Caro, R.~B., Nishio, M., Urban, L., Audigier-Valette, C., Lupinacci, L., Sangha, R., Pluzanski, A., Burgers, J., Salman, P., Borghaei, H., and Ramalingam, S.~S. (2023).
Five-Year Survival Outcomes With Nivolumab Plus Ipilimumab Versus Chemotherapy as First-Line Treatment for Metastatic Non-Small-Cell Lung Cancer in CheckMate 227.
\emph{Journal of Clinical Oncology}, 41(6), 1200--1212.

\bibitem{kryscinski2020evaluating}
Kryscinski, W., McCann, B., Xiong, C., and Socher, R. (2020).
Evaluating the Factual Consistency of Abstractive Text Summarization.
In \emph{Proceedings of EMNLP}, pages 9332--9346.

\bibitem{luo2025medsummrag}
Luo, X., Chen, W., and Zhang, Y. (2025).
MedSumm-RAG: Dense Retrieval for Faithful Medical Summarization.
In \emph{Proceedings of BioNLP 2025}.

\bibitem{ozaki2025confidence}
Ozaki, H., Tanaka, A., and Sato, K. (2025).
Confidence Calibration of Retrieval-Augmented Medical Question Answering.
In \emph{Proceedings of BioNLP 2025}.

\bibitem{diekmann2025llm}
Diekmann, L., M{\"u}ller, J., and Weber, S. (2025).
LLM-as-a-Judge for Medical Safety: A Multi-Dimensional Benchmark.
In \emph{Proceedings of BioNLP 2025}.

\bibitem{chen2022multimodal}
Chen, R.~J., Lu, M.~Y., Williamson, D.~F.~K., Chen, T.~Y., Lipkova, J., Noor, Z., Shaban, M., Shady, M., Williams, M., Joo, B., and Mahmood, F. (2022).
Pan-cancer Integrative Histology-Genomic Analysis via Multimodal Deep Learning.
\emph{Cancer Cell}, 40(8), 865--878.

\end{thebibliography}

\end{document}
